\documentclass[12pt]{article}
\usepackage[utf8]{inputenc}
\usepackage[T1]{fontenc}
\usepackage{lmodern}
\usepackage{microtype}
\usepackage[hidelinks]{hyperref}
\usepackage[a4paper, margin=3cm]{geometry}
\usepackage{parskip}

\title{What is Meditation}
\date{2026-08-02}

\begin{document}
\maketitle

Whenever I am at an event, where we invite people to try out Sahaja Yoga Meditation, I often get the answer: "No thanks, I already meditate."
And when I ask them which method they follow, I get answers ranging from: "I do mindfulness" to "When I am in the mountains it's like meditation for me."
Personally, I feel those are valid answers.
However, how can meditation be so ambigous?
Is it really the same if I do meditation according to the Buddhist tradition, mindfulness, Christian contemplation, Wim Hoff, or Sahaja Yoga?
I don't think so.

There seems to be a confusion about the English term \textit{meditation}.
First of all, there is meditation in the sense of contemplation or reflection --- just like in Marcus Aurelius' \textit{meditations}.
Then there is 'active meditation', where people describe a state of deep concentration, like the flow state.
If we go more toward the Indian tradition, meditation is a translation of \textit{dhyana} that is both used in Hindu and Buddhist tradition.
And ultimately there is meditation according to Sahaja Yoga that describes \textit{Nirvichara Samadhi}, or thoughtless awareness.
It seems as if \textit{meditation} is more of a conglomerate term rather than a specific activity.

To differentiate and cluster these different approaches I find Patanjali's framework useful.
In his \textit{Yoga Sutras} he describes 8 limbs of Yoga, where the first 4 talk more of physical or outward activites, while the final 4 are concerned with the inner being.
These are \textit{Pratyahara}, \textit{Dharana}, \textit{Dhyana}, and \textit{Samadhi}.
They describe the movement of the attention inside, the focusing of attention, then the cessation of thought, and finally the union with the self and the divine.

Now, if we apply that framework to some of the examples of different types of meditation from above, we can see that reflecting on our thoughts or behaviours is a way of moving the attention inside (\textit{Pratyahara}). Marcus Aurelius' famous book is the result of a kind of journaling, where he tried to look into his being to find knowledge.
An athlete, creative, or scientist, on the other hand, will describe the focused work, or deep work, as an activity that creates stillness, purpose, and connection with oneself.
Patanjali would call this \textit{Dharana}, which is an important step towards Yoga.
Also Mindfulness practitioners emphasize this type of meditation in their philosophy.
For Buddhists, or Hindus, however, that wouldn't count as meditation yet.
In their tradition meditation means \textit{Dhyana}, where the practitioner experiences the complete cessation of thought.
This thoughtlessness is created by the elongation of the gap between the thoughts and is the key to separating the attention from the mind.
When thoughts cease, the mind doesn't define the being anymore and a cleansing of the attention from the biases of the mind takes place.
Finally, the Sahaja Yoga meditation takes it one step further and corresponds to \textit{Samadhi}.
In that state the mind is transcended entirely, and the identification with the inner self is established.
It describes the true union, i.e. Yoga, and also acknowledges a stationary positioning in this state of awareness.
In Sahaja Yoga that is called \textit{thoughtless awareness} --- a state where the identification with inner self stays intact throughout and allows the entire being to reflect the light of spirit through every action or thought.
That is also the description of an enlightened, or self-realized being and requires the awakening of the Kundalini energy.

As we can see, there are multiple facettes to meditation making the term rather misleading.
Depending on the interpration, it can have very different meanings, where the two ends of the spectrum barely relate to each other.
However, using Patanjali's framework we can see how the relationship between the different methods.
\textit{Meditation} is used for any of the 4 final limbs of Yoga and makes it seem like the different types of mediations are uncorrelated, while they actually build on top of each other to reach the goal of Yoga, or enlightenment.

\end{document}
